\documentclass[a4paper,11pt]{article}
\usepackage{graphicx}
\usepackage[T1]{fontenc}
\usepackage[left=25mm, right=25mm, top=15mm, bottom=20mm, noheadfoot]{geometry}
\pagenumbering{gobble}
\usepackage{amsmath}

\usepackage{hyperref}
\usepackage{eso-pic, rotating, graphicx}

\AddToShipoutPicture{\put(30,200){\rotatebox{90}{\scalebox{1}{ \href{https://aca2026kosovo.net}{Applications of Computer Algebra (ACA) 2026 Conference.} }}}}

\begin{document}

% Please change the following abstract title and author description.
\title{Diagonalization and Change of Basis  in Linear Algebra}
\author{
 \underline{D. J. Jeffrey}$^{1*}$ \\
$^{1}$ University of Western Ontario \\
 email. djeffrey@uwo.ca
}

% Please do not modify following lines
\date{}
\maketitle

\noindent \textbf{Key words:}
\emph{%
% List your keywords here:
Linear Algebra; 
Similarity;
Basis change
}
\\\\
% Beginning of your abstract
A common topic in a first course on Linear Algebra is
diagonalization using eigenvectors \cite{poole}. Typically the textbooks define
``similarity'' as $B=V^{-1} A V$  with the justification that the eigenvalues
of $A$ and $B$ are the same. Later in the course (and in the text book)
there is a brief section on change of basis. No connection is made
between the topics. In this talk, I show the connection. It gives a
geometrical justification for the similarity transform, and some nice
plots. The change of basis connection extends to other diagonalizations
met in more advanced courses: topics such as the SVD factors of a matrix.

%You can cite articles or books such as \cite{} or \cite{}.


% End of your abstract

% Please add references below. 
\begin{thebibliography}{01}
\bibitem{poole} David Poole, Linear Algebra, 4th ed., ISBN 978-1-285-46324-7, Cengage Learning 2015.
\end{thebibliography}

\end{document}
